% Source: https://artofproblemsolving.com/community/c3h1311679_the_problem_that_has_baffled_entire_mathematics_departments
% Stable identity: aops-topic-1311679-q01
Consider triangle $DEC$.

\begin{asy}
usepackage("amsmath");
size(6cm);
defaultpen(fontsize(10pt));
pair D=(0,0),C=(5,0),A=(1.2,2.5),B=A+C; pair E=(3.0,2.5);
pair v=C-0.4*E; pair F=E+v, G=v;
draw(A--B--C--D--cycle,linewidth(0.9)); draw(D--E--F--G--cycle,blue);
draw(D--E--C--cycle,red+dashed);
label("$A$",A,NW); label("$B$",B,NE); label("$C$",C,E); label("$D$",D,W); label("$E$",E,N); label("$F$",F,SE,blue); label("$G$",G,SW,blue);
label("\mbox{}",(min(currentpicture,user=true)+max(currentpicture,user=true))/2,black); // keeps the text after the figure black when LaTeX typesets the labels inline
\end{asy}

Its base $DC$ is a side of $ABCD$, and its third vertex $E$ lies on the opposite side $AB$, so its height equals the height of $ABCD$: $[DEC]=\dfrac{1}{2}[ABCD]$.

Also, its base $DE$ is a side of $EFGD$, and its third vertex $C$ lies on the opposite side $FG$, so $[DEC]=\dfrac{1}{2}[EFGD]$.

Hence $[EFGD]=[ABCD]=\boxed{20}$.
